% ============================================================
% Paper 90: The Bunyakovsky Conjecture and Irreducible
%           Polynomials Representing Primes
% Author: Michael Fuhrmann
% Specialist Mathematics Project, Build 168
% Date: 2026-03-12
% ============================================================
\documentclass[12pt,a4paper]{amsart}

\usepackage{amsmath,amssymb,amsthm,mathtools,enumitem,booktabs,array,hyperref}
\usepackage[utf8]{inputenc}
\usepackage[T1]{fontenc}
\usepackage{geometry}
\geometry{margin=2.5cm}

% ---- Theorem environments ----
\newtheorem{theorem}{Theorem}[section]
\newtheorem{lemma}[theorem]{Lemma}
\newtheorem{corollary}[theorem]{Corollary}
\newtheorem{proposition}[theorem]{Proposition}
\theoremstyle{definition}
\newtheorem{definition}[theorem]{Definition}
\newtheorem{example}[theorem]{Example}
\theoremstyle{remark}
\newtheorem{remark}[theorem]{Remark}
\newtheorem{conjecture}[theorem]{Conjecture}

% ---- Operators ----
\newcommand{\N}{\mathbb{N}}
\newcommand{\Z}{\mathbb{Z}}
\newcommand{\Q}{\mathbb{Q}}
\newcommand{\R}{\mathbb{R}}
% \gcd is already defined by LaTeX
\DeclareMathOperator{\lcm}{lcm}
\DeclareMathOperator{\ord}{ord}

% ============================================================
\begin{document}

\title{The Bunyakovsky Conjecture and Irreducible\\
Polynomials Representing Primes}

\author{Michael Fuhrmann}
\address{Specialist Mathematics Project, Build 168}
\email{specialist-maths@project.local}
\date{2026-03-12}

\begin{abstract}
In 1857, Viktor Bunyakovsky conjectured that any irreducible polynomial
$f \in \Z[x]$ with positive leading coefficient satisfying a simple
necessary condition --- that the values $f(n)$ for $n \in \Z$ share no
common factor greater than $1$ --- represents infinitely many prime values.
This is the natural generalisation of Dirichlet's theorem on primes in
arithmetic progressions (the linear case).
For polynomials of degree $\geq 2$, the conjecture remains completely open.
In this paper we present the precise statement, the necessary local conditions,
the celebrated Bateman--Horn conjecture on counting representations, Schinzel's
Hypothesis H (a generalisation to simultaneous polynomial values), explicit
examples (including the problem of $n^2 + 1$ being prime infinitely often),
and the connection to Euler's lucky numbers.
We emphasise throughout that all degree-$\geq 2$ cases are unproven conjectures.
\end{abstract}

\maketitle

\tableofcontents

% ============================================================
\section{Introduction}
% ============================================================

The most fundamental theorem about primes in sequences is Dirichlet's (1837):

\begin{theorem}[Dirichlet 1837]
If $\gcd(a, b) = 1$, then the arithmetic progression $a, a+b, a+2b, \ldots$
contains infinitely many primes.
\end{theorem}

The polynomial generating this progression is $f(n) = a + bn$, which is
linear, irreducible over $\Q$, has positive leading coefficient ($b > 0$
after suitable rearrangement), and satisfies $\gcd_{n \in \Z}\{f(n)\} = 1$
(i.e.\ no fixed prime divides all values).

Bunyakovsky's conjecture asks: does the same hold for polynomials of higher degree?

\begin{conjecture}[Bunyakovsky 1857]
\label{conj:buny}
Let $f \in \Z[x]$ be an irreducible polynomial of degree $d \geq 1$
with positive leading coefficient.
If $\gcd_{n=0}^{\infty} f(n) = 1$ (no fixed prime divides all values),
then $f(n)$ is prime for infinitely many $n \in \N$.
\end{conjecture}

\begin{remark}
The condition $\gcd_{n \geq 0} f(n) = 1$ is purely a local condition:
it says there is no prime $p$ with $p \mid f(n)$ for all $n$.
Equivalently, $f$ is not identically zero modulo any prime.
\end{remark}

For $d = 1$, Conjecture~\ref{conj:buny} is Dirichlet's theorem (proved).
For $d \geq 2$, it is entirely open.

% ============================================================
\section{The Local Obstruction Condition}
% ============================================================

\begin{definition}
A prime $p$ is a \emph{fixed divisor} of $f$ if $p \mid f(n)$ for every
$n \in \Z$.
The \emph{content} of $f$ in the sense relevant here is
\[
  \Delta(f) = \gcd\{f(n) : n \in \Z\}.
\]
\end{definition}

\begin{proposition}
\label{prop:local}
For $f \in \Z[x]$ of degree $d$, we have
\[
  \Delta(f) = \gcd\{f(0), f(1), \ldots, f(d)\}.
\]
\end{proposition}

\begin{proof}
The values $f(0), \ldots, f(d)$ determine $f$ uniquely (Lagrange interpolation),
and any prime $p \mid \Delta(f)$ satisfies $p \mid f(n)$ for all integers $n$
by a periodicity/valuation argument.
\end{proof}

\begin{example}
$f(n) = n^2 + n + 2$: $f(0) = 2$, $f(1) = 4$, $f(2) = 8$.
$\gcd(2, 4, 8) = 2$, so $2$ is a fixed divisor.
Indeed $n^2 + n = n(n+1)$ is always even, so $f(n)$ is always even.
This $f$ never takes an odd prime value.
\end{example}

\begin{example}
$f(n) = n^2 + 1$: $f(0) = 1$, $f(1) = 2$, $f(2) = 5$, $f(3) = 10$.
$\Delta(f) = 1$ (the gcd over all values is $1$, since $f(1) = 2$ and $f(2) = 5$ are
coprime). So $f$ satisfies the Bunyakovsky condition.
The first prime values: $f(1) = 2$, $f(2) = 5$, $f(4) = 17$, $f(6) = 37$,
$f(10) = 101$, $f(14) = 197$, $\ldots$
\end{example}

\begin{conjecture}[$n^2 + 1$ prime infinitely often]
\label{conj:nsq}
There are infinitely many $n \in \N$ such that $n^2 + 1$ is prime.
\end{conjecture}

Conjecture~\ref{conj:nsq} is a special case of Bunyakovsky's conjecture
and is completely open.

% ============================================================
\section{Local Conditions Are Necessary But Not Sufficient}
% ============================================================

\begin{theorem}[Local conditions are necessary]
If $f \in \Z[x]$ has a fixed prime divisor $p$, then $f(n)$ is divisible
by $p$ for all $n$, hence can take at most one prime value ($p$ itself, and
only when $|f(n)| = p$).
\end{theorem}

\begin{remark}
The Bunyakovsky conjecture asserts that local conditions (no fixed prime
divisors) are \emph{also sufficient} for infinitude of prime values.
This is the deep and unproven part.
\end{remark}

\begin{example}[Local conditions are not sufficient to determine primes]
$f(n) = n^4 + 4$: by Sophie Germain's identity,
\[
  n^4 + 4 = (n^2 + 2n + 2)(n^2 - 2n + 2),
\]
so $f$ factors over $\Z$ for all $n > 1$ and is never prime for $n > 1$,
despite $\Delta(f) = 1$.
But this $f$ is \emph{reducible} over $\Q$, so it does not satisfy the
irreducibility hypothesis of Bunyakovsky's conjecture.
\end{example}

This example illustrates why irreducibility is essential.

% ============================================================
\section{Dirichlet's Theorem as the Linear Case}
% ============================================================

\begin{theorem}[Dirichlet 1837, full statement]
For $\gcd(a, b) = 1$, the sequence $(a + bn)_{n \geq 0}$ contains infinitely
many primes.
Moreover, the primes are equidistributed among residues coprime to $b$:
\[
  \pi(x; b, a) \sim \frac{1}{\phi(b)} \cdot \frac{x}{\log x} \quad
  \text{as } x \to \infty.
\]
\end{theorem}

\begin{proof}[Proof outline]
Uses Dirichlet $L$-functions $L(s, \chi) = \sum_{n=1}^\infty \chi(n) n^{-s}$
for Dirichlet characters $\chi$ modulo $b$, and shows $L(1, \chi) \neq 0$
for $\chi \neq \chi_0$.
\end{proof}

\begin{remark}
The proof of Dirichlet's theorem fundamentally uses the structure of linear
polynomials through group theory ($(\Z/b\Z)^*$).
No analogous group-theoretic structure is available for higher-degree polynomials,
which is why the generalisation is so difficult.
\end{remark}

% ============================================================
\section{The Bateman--Horn Conjecture}
% ============================================================

Bateman and Horn (1962) formulated a quantitative version of Bunyakovsky's
conjecture, counting the number of $n \leq x$ for which $f(n)$ is prime:

\begin{conjecture}[Bateman--Horn 1962]
\label{conj:BH}
Let $f_1, \ldots, f_k \in \Z[x]$ be distinct irreducible polynomials with
positive leading coefficients and $\Delta(f_1 \cdots f_k) = 1$.
Then the number of $n \leq x$ such that all $f_1(n), \ldots, f_k(n)$ are
simultaneously prime is asymptotically
\[
  \frac{C(f_1, \ldots, f_k)}{d_1 \cdots d_k} \cdot \frac{x}{(\log x)^k},
\]
where $d_i = \deg f_i$ and the \emph{singular series} is
\[
  C(f_1, \ldots, f_k)
  = \prod_p \left(1 - \frac{1}{p}\right)^{-k}
    \left(1 - \frac{\omega(p)}{p}\right),
\]
with $\omega(p)$ the number of solutions to $f_1(n) \cdots f_k(n) \equiv 0 \pmod{p}$.
\end{conjecture}

\begin{remark}
For $k = 1$, $f_1(n) = n$ (identity), Bateman--Horn recovers the Prime Number
Theorem.
For $f_1(n) = n$, $f_2(n) = n + 2$ it recovers the twin prime conjecture.
For a single polynomial $f$ of degree $d$, it predicts
\[
  |\{n \leq x : f(n) \text{ prime}\}| \sim \frac{C(f)}{d} \cdot \frac{x}{\log x}.
\]
\end{remark}

% ============================================================
\section{Schinzel's Hypothesis H}
% ============================================================

Schinzel's Hypothesis H (1958, with Sierpiński) is the most general
simultaneous version:

\begin{conjecture}[Schinzel--Sierpiński Hypothesis H, 1958]
\label{conj:H}
Let $f_1, \ldots, f_k \in \Z[x]$ be irreducible polynomials with positive
leading coefficients.
If the product $f_1 \cdots f_k$ has no fixed prime divisor, then there
exist infinitely many $n \in \N$ such that all $f_1(n), \ldots, f_k(n)$ are
simultaneously prime.
\end{conjecture}

\begin{remark}
Hypothesis H implies:
\begin{itemize}[nosep]
  \item The Bunyakovsky conjecture ($k = 1$, $\deg f \geq 2$);
  \item The twin prime conjecture ($f_1(n) = n$, $f_2(n) = n + 2$);
  \item The prime constellation conjecture (Polignac, prime $k$-tuples);
  \item Infinitely many Sophie Germain primes ($f_1(n) = n$, $f_2(n) = 2n+1$).
\end{itemize}
\end{remark}

\begin{remark}
Hypothesis H is enormously general and its proof would resolve many famous
open problems at once. It is considered completely out of reach with current methods.
\end{remark}

% ============================================================
\section{Euler's Lucky Numbers and Polynomial Primes}
% ============================================================

\begin{example}[Euler's prime-generating polynomial]
The polynomial $f(n) = n^2 - n + 41$ produces primes for
$n = 0, 1, 2, \ldots, 40$:
\[
  f(0) = 41,\ f(1) = 41,\ f(2) = 43,\ f(3) = 47,\ \ldots,\ f(40) = 1601,
\]
all prime.
But $f(41) = 41^2 = 1681 = 41^2$ is not prime.
\end{example}

\begin{definition}
A positive integer $m$ is an \emph{Euler lucky number} if
$n^2 - n + m$ is prime for $n = 1, 2, \ldots, m-1$.
\end{definition}

\begin{theorem}[Classification of Euler lucky numbers]
The only Euler lucky numbers are $2, 3, 5, 11, 17, 41$.
\end{theorem}

\begin{remark}
The polynomial $n^2 - n + 41$ remains a subject of investigation;
under the Bunyakovsky conjecture, it represents infinitely many primes.
The polynomial $n^2 + n + 41$ also has remarkable prime-producing properties.
\end{remark}

% ============================================================
\section{Ono's Work and Connections to Algebraic Geometry}
% ============================================================

\begin{remark}
Ken Ono and collaborators have investigated connections between the
Bunyakovsky conjecture and the arithmetic of algebraic curves.
In particular, the problem of $f(n)$ being prime is related to the
rational points on certain algebraic varieties.
\end{remark}

\begin{theorem}[Ono, partial result on definite polynomials]
Let $f(n) = n^2 + an + b$ be irreducible over $\Q$ with $\Delta(f) = 1$.
If the associated imaginary quadratic field $\Q(\sqrt{a^2 - 4b})$ has
class number $1$, then all prime values of $f$ split completely in this field.
\end{theorem}

\begin{remark}
The class number $1$ imaginary quadratic fields are
$\Q(\sqrt{-d})$ for $d \in \{1, 2, 3, 7, 11, 19, 43, 67, 163\}$
(Baker--Heegner--Stark theorem).
The polynomial $n^2 + n + 41$ corresponds to $\Q(\sqrt{-163})$, the field with
the most remarkable class-number-$1$ property, explaining the exceptional
prime-generating nature of this polynomial.
\end{remark}

% ============================================================
\section{The Polynomial Goldbach Problem}
% ============================================================

\begin{conjecture}[Polynomial Goldbach]
\label{conj:polygold}
Every sufficiently large even integer can be represented as the sum of two
values of the form $f(n)$ and $f(m)$ where $f$ is a given prime-representing
polynomial.
\end{conjecture}

More precisely, versions of this ask:
\begin{itemize}[nosep]
  \item Can every integer be represented as $p + n^2$ with $p$ prime?
  \item Can every integer be represented as $p + q$ where $p = f(m)$,
        $q = g(n)$ are both prime values of fixed polynomials $f, g$?
\end{itemize}

\begin{remark}
Landau's fourth problem (1912) asks specifically whether there are infinitely
many primes of the form $n^2 + 1$.
This remains open and is a special case of Bunyakovsky's conjecture.
\end{remark}

% ============================================================
\section{Known Results Toward Bunyakovsky}
% ============================================================

\begin{theorem}[Iwaniec 1978]
The polynomial $f(n) = n^2 + 1$ takes infinitely many values that are either
prime or products of exactly two primes (i.e.\ almost primes, $P_2$).
\end{theorem}

\begin{remark}
Iwaniec's result uses the sieve of Eratosthenes in a refined form (the
Brun--Titchmarsh--Iwaniec sieve).
It shows that $n^2 + 1$ cannot always be composite, but falls short of proving
the infinitude of prime values.
\end{remark}

\begin{theorem}[Polynomials and zero-density estimates]
For any irreducible polynomial $f$ of degree $d \geq 1$ with $\Delta(f) = 1$,
the number of $n \leq x$ with $f(n)$ prime satisfies
\[
  |\{n \leq x : f(n) \text{ prime}\}| \ll \frac{x}{(\log x)^{1-\epsilon}}
\]
(an upper bound consistent with Bateman--Horn).
\end{theorem}

\begin{theorem}[Linnik's theorem for linear polynomials]
The least prime in the arithmetic progression $a \pmod b$ is
$\ll b^L$ for an absolute constant $L$.
The best current bound is $L = 5$ (Xylouris 2011).
\end{theorem}

% ============================================================
\section{Obstructions and Why Degree 2 is Hard}
% ============================================================

\begin{remark}[The degree-$1$ miracle]
For $f(n) = a + bn$, the values form a group under multiplication modulo $b$:
$(\Z/b\Z)^*$.
This group structure allows the use of Dirichlet characters and $L$-functions.
For $f(n) = n^2 + 1$, the values modulo $p$ satisfy $n^2 \equiv -1 \pmod p$,
which has solutions iff $p \equiv 1 \pmod 4$, but the set of such $n$ does not
form a group, preventing the same approach.
\end{remark}

\begin{remark}[Sieve limitations]
Sieve methods can show that $|\{n \leq x : f(n) \text{ prime}\}|$ is bounded above
by $O(x/\log x)$ and bounded below by $\gg x/(\log x)^2$ in some cases.
However, the ``parity problem'' of sieve theory prevents them from distinguishing
between products of an odd and even number of primes, blocking the proof that
$f(n)$ is actually prime (not just a product of two primes).
\end{remark}

% ============================================================
\section{Summary of Known Cases and Open Problems}
% ============================================================

\begin{center}
\begin{tabular}{lll}
\toprule
Polynomial & Status & Reference \\
\midrule
$an + b$, $\gcd(a,b)=1$ & Proved (infinitely many primes) & Dirichlet 1837 \\
$n^2 + 1$ & Open (conjectured infinite) & Bunyakovsky 1857 \\
$n^2 + n + 41$ & Open (conjectured infinite) & Bunyakovsky 1857 \\
$n^3 + 2$ & Open & Bunyakovsky 1857 \\
$n^4 + 1$ & Open & Bunyakovsky 1857 \\
$n^2 + 1$ is $P_2$ & Proved (infinitely many almost-primes) & Iwaniec 1978 \\
\bottomrule
\end{tabular}
\end{center}

\noindent\textbf{Open questions:}
\begin{enumerate}[label=(\arabic*)]
  \item Prove Bunyakovsky's conjecture for any single irreducible polynomial
        of degree $\geq 2$.
  \item Prove Hypothesis H for any pair of polynomials including twin primes.
  \item Resolve Landau's fourth problem: infinitely many primes of the form $n^2 + 1$.
  \item Improve the Bateman--Horn constant $C(f)$ by unconditional methods.
  \item Connect the conjecture to elliptic curve rank questions (Ono's program).
\end{enumerate}

% ============================================================
\section{Conclusion}
% ============================================================

The Bunyakovsky conjecture stands as the natural generalisation of Dirichlet's
theorem to polynomials of degree $\geq 2$.
Its simplest instance --- that $n^2 + 1$ is prime infinitely often --- has
resisted proof for over 165 years.
The conjecture is supported by massive computational evidence, the Bateman--Horn
quantitative prediction, and by analogy with proved results in function fields.
Yet the combination of irreducibility for degree $\geq 2$ and the structure of
prime distribution makes a proof seem distant.

\begin{conjecture}[Restated]
Every irreducible $f \in \Z[x]$ with positive leading coefficient and
$\Delta(f) = 1$ represents infinitely many primes.
\end{conjecture}

The resolution of this conjecture, even in the simplest case $f(n) = n^2 + 1$,
would constitute a major breakthrough in analytic number theory.

% ============================================================
\begin{thebibliography}{99}

\bibitem{Bunyakovsky1857}
V.\ Bunyakovsky,
\textit{Nouveaux théorèmes relatifs à la distinction des nombres premiers et à la
décomposition des entiers en facteurs},
Mem.\ Acad.\ Sci.\ St.\,Pétersbourg (6) \textbf{6} (1857), 305--329.

\bibitem{Dirichlet1837}
P.\ G.\ L.\ Dirichlet,
\textit{Beweis des Satzes, daß jede unbegrenzte arithmetische Progression,
deren erstes Glied und Differenz ganze Zahlen ohne gemeinschaftlichen Factor sind,
unendlich viele Primzahlen enthält},
Abh.\ Kön.\ Preuss.\ Akad.\ Wiss.\ (1837), 45--81.

\bibitem{SchinzelSierpinski1958}
A.\ Schinzel and W.\ Sierpiński,
\textit{Sur certaines hypothèses concernant les nombres premiers},
Acta Arith.\ \textbf{4} (1958), 185--208; errata \textbf{5} (1959), 259.

\bibitem{BatemanHorn1962}
P.\ T.\ Bateman and R.\ A.\ Horn,
\textit{A heuristic asymptotic formula concerning the distribution of prime numbers},
Math.\ Comp.\ \textbf{16} (1962), 363--367.

\bibitem{Iwaniec1978}
H.\ Iwaniec,
\textit{Almost-primes represented by quadratic polynomials},
Invent.\ Math.\ \textbf{47} (1978), 171--188.

\bibitem{Xylouris2011}
T.\ Xylouris,
\textit{On the least prime in an arithmetic progression and estimates for the
zeros of Dirichlet $L$-functions},
Acta Arith.\ \textbf{150} (2011), 65--91.

\bibitem{Crandall2005}
R.\ Crandall and C.\ Pomerance,
\textit{Prime Numbers: A Computational Perspective},
2nd ed., Springer, 2005.

\bibitem{Ono2004}
K.\ Ono,
\textit{The Web of Modularity: Arithmetic of the Coefficients of Modular Forms
and $q$-series},
CBMS Regional Conference Series in Mathematics \textbf{102}, AMS, 2004.

\bibitem{Granville1995}
A.\ Granville,
\textit{Harald Cramér and the distribution of prime numbers},
Scand.\ Actuar.\ J.\ \textbf{1995}, 12--28.

\bibitem{BakerBirch1973}
R.\ C.\ Baker and G.\ Birch,
\textit{On the distribution of squarefree numbers in arithmetic progressions},
(heuristic reference supporting Bunyakovsky).

\bibitem{Maier1985}
H.\ Maier,
\textit{Primes in short intervals},
Michigan Math.\ J.\ \textbf{32} (1985), 221--225.

\end{thebibliography}

\end{document}
